\section{Method}

Motivated by the pivotal insight that positional information dominates query representations, we propose DapQ \footnote{\url{https://github.com/tianzhenxu/DapQ}}(as illustrated in Figure \ref{fig:DapQ}), a novel KV cache compression framework that accurately simulates decoding-stage contextual positioning during the prefill phase. DapQ synthesizes a decoding-aligned observation window, composed of pseudo queries endowed with future positional encodings, which mirrors the dynamic context of the actual decoding process. This precisely assesses token importance, enabling accurate token eviction without altering the intended timeline.

\subsection{Construct Decoding-Aligned \texorpdfstring{\(Q_{\text{pseudo}}\)}{Q\_pseudo}}
% \subsection{Construct Decoding-Aligned \(Q_{\text{pseudo}}\)}
The core of DapQ is to simulate the dynamic positional query of the decoding phase. For a prompt sequence of length \(L_p\), we append a set of \(N\) artificially constructed tokens, denoted \(\mathbf{T}_{\text{pseudo}}\), to form an extended input sequence. \(\mathbf{T}_{\text{pseudo}}\) can be constructed or arbitrarily chosen from the existing context (e.g., uniformly sampled or prefix-suffix concatenation), as their semantic content is secondary to the positional assignment. The crucial operation is to assign \(\mathbf{T}_{\text{pseudo}}\) the correct positional indices that they would occupy as the first \(N\) tokens generated by the model, rather than arbitrarily:
\begin{equation*}
\text{Positions}(\mathbf{T}_{\text{pseudo}}) = [L_p, L_p+1, ..., L_p+N-1].
\end{equation*}
This yields an input sequence with length \(L_{\text{total}} = L_p + N\). The model processes this extended sequence during the prefill phase, computing KV cache for \(L_{\text{p}}\) prompt tokens. The primary purpose of this step is to obtain pseudo queries (\(Q_{\text{pseudo}}\)) of these  \(\mathbf{T}_{\text{pseudo}}\), which are endowed with correct positional encodings for the start of decoding phase.

\subsection{Importance Assessment and Eviction}
We leverage the \(Q_{\text{pseudo}}\) to assess the importance of all Keys derived from the original prompt. The importance score for the \(j\)-th (\(j \in [0, L_p - 1]\)) prompt token is computed by aggregating its attention scores from each pseudo query \(q_i \in Q_{\text{pseudo}}\):
\begin{equation}
S(j) = \sum_{i=L_p}^{L_p+N-1} \text{Attention}(q_i, k_j).
\end{equation}
The \(TopK\) tokens with the highest scores \(S(j)\) are retained:
\begin{equation}
M_{\text{retain}} = \underset{j \in [0, L_p - 1]}{TopK}\left( S(j) \right).
\end{equation}
The KV cache is pruned, discarding all key-value pairs not in \(M_{\text{retain}}\). \textbf{Crucially, the entire synthetic segment \(\mathbf{T}_{\text{pseudo}}\) is discarded immediately after performing the importance scoring.} Autoregressive decoding phase then begins from position \(L_p\), utilizing only the compressed cache of size \(K\). This ensures the model's generation remains consistent with the intended timeline.



